% =====================================================================
\chapter{Soluciones de los ejercicios}
\label{ap:soluciones}
% =====================================================================

Cada solución se presenta en cinco partes: datos, planteamiento, desarrollo, resultado y comprobación. En los pasos intermedios se redondea a cuatro cifras significativas salvo que se indique otra cosa; los resultados exactos se dan como fracción cuando existe.

% ---------------------------------------------------------------------
\section{Tema 1}

\begin{solucion}{eje:1-decod}
\datos un decodificador que aplica la decodificación UTF-8 a cada byte por separado.
\planteamiento un carácter no ASCII ocupa varios bytes (\cref{tab:utf8}) y ninguno de ellos es por sí solo una palabra del código.
\pasos
\begin{enumerate}[label=\arabic*.]
  \item (a) «año» se codifica como $(97,195,177,111)=(\hx{61},\hx{C3},\hx{B1},\hx{6F})$. El byte $97$ da «a», pero $\hx{C3}=\bits{11000011}$ es un byte inicial que anuncia un carácter de 2 bytes; aislado, la cadena está incompleta y la decodificación falla.
  \item El error de fondo es suponer que un byte es un carácter, lo que solo es cierto en ASCII. Con «hola» el programa funciona, y por eso el fallo pasa inadvertido.
  \item (b) La cadena $\hx{C3}\,\hx{28}$: el primer byte exige un segundo byte que empiece por \bits{10}, pero $\hx{28}=\bits{00101000}$ empieza por \bits{00}. Otra respuesta válida es $\hx{FF}\,\hx{FF}$, porque ningún byte UTF-8 empieza por \bits{11111}.
\end{enumerate}
\resultado \textbf{(a) «año» falla porque la «ñ» $=\hx{C3 B1}$ se parte y ninguna mitad es un carácter; (b) $\hx{C3 28}$ (o $\hx{FF FF}$) no es UTF-8 válido.}
\comprobacion en Python, \texttt{b"\textbackslash xc3\textbackslash x28".decode()} lanza el error \emph{invalid continuation byte} y \texttt{b"\textbackslash xff\textbackslash xff".decode()}, \emph{invalid start byte}.
\end{solucion}

\begin{solucion}{eje:1-e}
\datos «é» $=\hx{U+00E9}$.
\planteamiento $\hx{E9}>\hx{7F}$, luego se usa la fila de 2 bytes, \bits{110xxxxx 10xxxxxx}, con 11 bits útiles.
\pasos
\begin{enumerate}[label=\arabic*.]
  \item Hexadecimal a binario: $\hx{E}=\bits{1110}$ y $\hx{9}=\bits{1001}$, luego $\hx{E9}=\bits{11101001}=233$.
  \item Con 11 bits: \bits{00011101001}; en grupos de $5+6$: \bits{00011} y \bits{101001}.
  \item Primer byte: $\bits{110}\,\bits{00011}=\bits{11000011}=\hx{C3}=195$.
  \item Segundo byte: $\bits{10}\,\bits{101001}=\bits{10101001}=\hx{A9}=169$.
\end{enumerate}
\resultado \textbf{$\bits{11000011}\ \bits{10101001}=\hx{C3 A9}$.}
\comprobacion al quitar las cabeceras \bits{110} y \bits{10} queda $\bits{00011}\,\bits{101001}=\hx{E9}$. En Python, \texttt{"é".encode()} da \texttt{b"\textbackslash xc3\textbackslash xa9"}.
\end{solucion}

\begin{solucion}{eje:1-contar}
\datos «pingüino», «Ñandú» y «café» seguido de espacio y de \hx{U+2615}.
\planteamiento sumar los bytes de cada carácter según su fila en el \cref{tab:utf8}.
\pasos
\begin{enumerate}[label=\arabic*.]
  \item «pingüino» tiene 8 caracteres: siete ASCII (p, i, n, g, i, n, o), que suman $7\cdot 1=7$ bytes, y «ü» (\hx{U+00FC}), con 2. Total: $7+2=9$.
  \item «Ñandú» tiene 5 caracteres: «Ñ» (\hx{U+00D1}) con 2, «a», «n», «d» con $3\cdot 1=3$ y «ú» (\hx{U+00FA}) con 2. Total: $2+3+2=7$.
  \item El tercer texto tiene 6 caracteres: «c», «a», «f» ($3$ bytes), «é» ($2$), el espacio ($1$) y \hx{U+2615}, que está en el intervalo \hx{U+0800}--\hx{U+FFFF} y ocupa $3$. Total: $3+2+1+3=9$.
\end{enumerate}
\resultado \textbf{9, 7 y 9 bytes.}
\comprobacion bytes menos caracteres: $9-8=1$ («ü»), $7-5=2$ («Ñ» y «ú») y $9-6=3$ («é» aporta 1 y \hx{U+2615} aporta 2).
\end{solucion}

\begin{solucion}{eje:1-mano}
\datos «¿» $=\hx{U+00BF}$; carácter chino \hx{U+8FDB}.
\planteamiento \hx{BF} está en la fila de 2 bytes (11 bits útiles); \hx{8FDB} está en la de 3 bytes (16 bits útiles, patrón \bits{1110xxxx 10xxxxxx 10xxxxxx}).
\pasos
\begin{enumerate}[label=\arabic*.]
  \item «¿»: $\hx{BF}=\bits{10111111}=191$. Con 11 bits, \bits{00010111111}; en grupos de $5+6$, \bits{00010} y \bits{111111}.
  \item Bytes: $\bits{110}\,\bits{00010}=\bits{11000010}=\hx{C2}$ y $\bits{10}\,\bits{111111}=\bits{10111111}=\hx{BF}$.
  \item \hx{U+8FDB}: $\hx{8FDB}=\bits{1000\,1111\,1101\,1011}$ (16 bits); en grupos de $4+6+6$, \bits{1000}, \bits{111111} y \bits{011011}.
  \item Bytes: $\bits{1110}\,\bits{1000}=\hx{E8}$, $\bits{10}\,\bits{111111}=\hx{BF}$ y $\bits{10}\,\bits{011011}=\bits{10011011}=\hx{9B}$.
\end{enumerate}
\resultado \textbf{«¿» $=\hx{C2 BF}$ (2 bytes); \hx{U+8FDB} $=\hx{E8 BF 9B}$ (3 bytes).}
\comprobacion al quitar las cabeceras del segundo carácter queda $\bits{1000}\,\bits{111111}\,\bits{011011}=\hx{8FDB}$. Obsérvese que «¿» empieza por \hx{C2} y no por \hx{C3}, como las vocales con tilde, porque $\hx{BF}<\hx{C0}$.
\end{solucion}

\begin{solucion}{eje:1-filas}
\datos el \cref{tab:utf8}; emoji \hx{U+1F600}.
\planteamiento cardinal de cada intervalo y cociente entre la fila de 4 bytes y el total de valores escalares.
\pasos
\begin{enumerate}[label=\arabic*.]
  \item 1 byte: $\hx{80}=128$. 2 bytes: $\hx{800}-\hx{80}=2048-128=1920$.
  \item 3 bytes: $\hx{10000}-\hx{800}=65\,536-2048=63\,488$, de los que $2048$ son sustitutos.
  \item 4 bytes: $\hx{110000}-\hx{10000}=1\,114\,112-65\,536=1\,048\,576=2^{20}$.
  \item Valores escalares: $128+1920+(63\,488-2048)+1\,048\,576=1\,112\,064$.
  \item Fracción con 4 bytes: $1\,048\,576/1\,112\,064=0{,}9429$.
  \item $\hx{1F600}\ge\hx{10000}$, luego está en la última fila. Sus bytes son $\hx{F0 9F 98 80}$.
\end{enumerate}
\resultado \textbf{el $94{,}29\,\%$ de los valores escalares necesita 4 bytes}; el emoji ocupa 4 bytes porque su punto de código supera \hx{FFFF}.
\comprobacion $128+1920+63\,488+1\,048\,576=1\,114\,112=\hx{110000}$, el número total de puntos de código. Casi todos los caracteres de uso frecuente están, sin embargo, en las tres primeras filas.
\end{solucion}

\begin{solucion}{eje:1-sobrelargo}
\datos las filas del \cref{tab:utf8} sin la condición de usar la fila más corta.
\planteamiento (a) exhibir dos cadenas de bytes distintas que se decodifiquen como el mismo carácter; (b) calcular el valor máximo que admite un byte inicial \hx{C0} o \hx{C1}.
\pasos
\begin{enumerate}[label=\arabic*.]
  \item (a) «a» $=\hx{61}$ tiene, con 11 bits, la escritura \bits{00001100001}. En la fila de 2 bytes daría $\bits{110}\,\bits{00001}=\hx{C1}$ y $\bits{10}\,\bits{100001}=\hx{A1}$.
  \item Las cadenas $(\hx{61})$ y $(\hx{C1},\hx{A1})$ son distintas y se decodificarían ambas como «a»: la decodificación no sería inyectiva.
  \item Para un tokenizador de bytes, un mismo texto tendría varias representaciones con tokens distintos, y el modelo vería como diferentes dos entradas iguales. Además, un filtro que busque un carácter por sus bytes podría ser burlado con la forma sobrelarga.
  \item (b) \hx{C0} $=\bits{11000000}$ y \hx{C1} $=\bits{11000001}$ aportan los 5 bits útiles \bits{0000x}. El valor de 11 bits resultante es como mucho $\bits{00001111111}=127$.
  \item Todo valor $\le 127$ cabe en la fila de 1 byte, luego una cadena que empiece por \hx{C0} o \hx{C1} es necesariamente sobrelarga.
\end{enumerate}
\resultado \textbf{(a) $(\hx{61})$ y $(\hx{C1},\hx{A1})$ representarían ambas «a»; (b) \hx{C0} y \hx{C1} solo codifican valores menores que~128.}
\comprobacion el byte inicial más pequeño válido de la fila de 2 bytes es \hx{C2}, que corresponde a $\hx{80}=128$, el primer punto de código de esa fila.
\end{solucion}

% ---------------------------------------------------------------------
\section{Tema 2}

\begin{solucion}{eje:2-unk}
\datos BPE a nivel de byte; $\V_0=\mathcal{B}$.
\planteamiento basta ver que todo texto se escribe con los símbolos de $\V_0$ y que las fusiones no los eliminan.
\pasos
\begin{enumerate}[label=\arabic*.]
  \item Todo texto, en cualquier idioma, se codifica en UTF-8 como una cadena de bytes (\cref{cor:vocab256}).
  \item Un byte tiene 8 bits y toma $2^{8}=256$ valores, todos en $\V_0$.
  \item Las fusiones solo añaden símbolos (\cref{prop:tam-vocab}); en el peor caso no se aplica ninguna y cada byte es un token (\cref{cor:no-unk}).
\end{enumerate}
\resultado \textbf{no hay tokens desconocidos porque los 256 bytes están en el vocabulario; el vocabulario inicial mide 256.}
\comprobacion la frase rusa de la \cref{tema:3}, con 22 bytes, da 22 tokens con el tokenizador de la práctica: ninguna fusión aplica y aun así se codifica.
\end{solucion}

\begin{solucion}{eje:2-fusiones}
\datos (a) $n$ símbolos iniciales y objetivo $m$; (b) $256$ bytes, $1$ token especial y $K=50\,257$.
\planteamiento cada fusión añade un símbolo (\cref{prop:tam-vocab}).
\pasos
\begin{enumerate}[label=\arabic*.]
  \item (a) $n+k=m$, luego $k=m-n$.
  \item (b) El vocabulario que no procede de fusiones tiene $256+1=257$ símbolos, luego $k=50\,257-257=50\,000$.
\end{enumerate}
\resultado \textbf{(a) $m-n$ fusiones; (b) $50\,000$ fusiones.}
\comprobacion $256+50\,000+1=50\,257$. El número redondo indica que $k=50\,000$ se fijó como hiperparámetro.
\end{solucion}

\begin{solucion}{eje:2-rana}
\datos «rana» $\times1$, «lana» $\times2$, «ana» $\times1$, «plan» $\times1$; 3 fusiones.
\planteamiento contar pares ponderados por la frecuencia, fusionar el máximo y restar su frecuencia (\cref{prop:compresion}). No hay empates en cabeza.
\pasos
\begin{enumerate}[label=\arabic*.]
  \item Símbolos iniciales: $4\cdot 1+4\cdot 2+3\cdot 1+4\cdot 1=4+8+3+4=19$.
  \item Paso 1: $N(\tok{a},\tok{n})=1+2+1+1=5$; $N(\tok{n},\tok{a})=1+2+1=4$; $N(\tok{l},\tok{a})=2+1=3$; $N(\tok{r},\tok{a})=N(\tok{p},\tok{l})=1$. Se fusiona $(\tok{a},\tok{n})$: \tok{r}\sep\tok{an}\sep\tok{a}; \tok{l}\sep\tok{an}\sep\tok{a}$\times2$; \tok{an}\sep\tok{a}; \tok{p}\sep\tok{l}\sep\tok{an}, con $19-5=14$ símbolos.
  \item Paso 2: $N(\tok{an},\tok{a})=1+2+1=4$; $N(\tok{l},\tok{an})=2+1=3$; $N(\tok{r},\tok{an})=N(\tok{p},\tok{l})=1$. Se fusiona $(\tok{an},\tok{a})$: \tok{r}\sep\tok{ana}; \tok{l}\sep\tok{ana}$\times2$; \tok{ana}; \tok{p}\sep\tok{l}\sep\tok{an}, con $14-4=10$ símbolos.
  \item Paso 3: $N(\tok{l},\tok{ana})=2$; $N(\tok{r},\tok{ana})=N(\tok{p},\tok{l})=N(\tok{l},\tok{an})=1$. Se fusiona $(\tok{l},\tok{ana})$: \tok{r}\sep\tok{ana}; \tok{lana}$\times2$; \tok{ana}; \tok{p}\sep\tok{l}\sep\tok{an}, con $10-2=8$ símbolos.
\end{enumerate}
\resultado \textbf{fusiones $\tok{a}+\tok{n}$ (5), $\tok{an}+\tok{a}$ (4) y $\tok{l}+\tok{ana}$ (2); símbolos $19\to 14\to 10\to 8$.}
\comprobacion recuento directo del corpus final: $2+1\cdot 2+1+3=8$.
\end{solucion}

\begin{solucion}{eje:2-codificar}
\datos fusiones en orden: (1) $\tok{a}+\tok{s}$, (2) $\tok{c}+\tok{as}$, (3) $\tok{cas}+\tok{a}$, (4) $\tok{cas}+\tok{as}$.
\planteamiento aplicar las fusiones una a una, en el orden aprendido (\cref{def:cod-bpe}).
\pasos
\begin{enumerate}[label=\arabic*.]
  \item «masas» $=\tok{m}\sep\tok{a}\sep\tok{s}\sep\tok{a}\sep\tok{s}$. (1) da $\tok{m}\sep\tok{as}\sep\tok{as}$; (2) no aplica porque no hay \tok{c}; (3) y (4) no aplican porque no hay \tok{cas}. Resultado: $\tok{m}\,\sep{}\,\tok{as}\,\sep{}\,\tok{as}$.
  \item «casa» $=\tok{c}\sep\tok{a}\sep\tok{s}\sep\tok{a}$. (1) da $\tok{c}\sep\tok{as}\sep\tok{a}$; (2), $\tok{cas}\sep\tok{a}$; (3), $\tok{casa}$; (4) no aplica. Resultado: un token.
  \item «cascas» $=\tok{c}\sep\tok{a}\sep\tok{s}\sep\tok{c}\sep\tok{a}\sep\tok{s}$. (1) da $\tok{c}\sep\tok{as}\sep\tok{c}\sep\tok{as}$; (2), $\tok{cas}\sep\tok{cas}$; (3) no aplica porque tras \tok{cas} viene \tok{cas} y no \tok{a}; (4) tampoco, porque tras \tok{cas} no viene \tok{as}. Resultado: $\tok{cas}\,\sep{}\,\tok{cas}$.
\end{enumerate}
\resultado \textbf{«masas» $\to\tok{m}\sep{}\tok{as}\sep{}\tok{as}$ (3 tokens); «casa» $\to\tok{casa}$ (1); «cascas» $\to\tok{cas}\sep{}\tok{cas}$ (2).}
\comprobacion al concatenar se recupera cada palabra, como garantiza el \cref{teo:bpe-correcto}.
\end{solucion}

\begin{solucion}{eje:2-lex}
\datos «paso» $\times4$, «pasos» $\times2$, «piso» $\times3$, «pisos» $\times3$; desempate por el par lexicográficamente mayor.
\planteamiento como en el \cref{ej:bpe-casa}, aplicando la regla de desempate cuando haga falta.
\pasos
\begin{enumerate}[label=\arabic*.]
  \item Símbolos iniciales: $4\cdot 4+2\cdot 5+3\cdot 4+3\cdot 5=16+10+12+15=53$.
  \item Paso 1: $N(\tok{s},\tok{o})=4+2+3+3=12$; $N(\tok{p},\tok{a})=N(\tok{a},\tok{s})=N(\tok{p},\tok{i})=N(\tok{i},\tok{s})=6$; $N(\tok{o},\tok{s})=2+3=5$. Sin empate: se fusiona $(\tok{s},\tok{o})$ y quedan $53-12=41$ símbolos.
  \item Paso 2: $(\tok{p},\tok{a})$, $(\tok{a},\tok{so})$, $(\tok{p},\tok{i})$ e $(\tok{i},\tok{so})$ empatan a 6; $(\tok{so},\tok{s})$ tiene 5. El primer símbolo mayor es \tok{p}; entre $(\tok{p},\tok{a})$ y $(\tok{p},\tok{i})$ gana \tok{i}. Se fusiona $(\tok{p},\tok{i})$: $41-6=35$.
  \item Paso 3: $(\tok{p},\tok{a})$, $(\tok{a},\tok{so})$ y $(\tok{pi},\tok{so})$ empatan a 6. Como $\tok{pi}>\tok{p}>\tok{a}$, se fusiona $(\tok{pi},\tok{so})$: $35-6=29$.
  \item Paso 4: $(\tok{p},\tok{a})$ y $(\tok{a},\tok{so})$ empatan a 6; $(\tok{piso},\tok{s})$ tiene 3 y $(\tok{so},\tok{s})$, 2. Gana $(\tok{p},\tok{a})$ porque $\tok{p}>\tok{a}$: $29-6=23$.
  \item (b) Fusiones en orden: \tok{so}, \tok{pi}, \tok{piso}, \tok{pa}. «peso» $=\tok{p}\sep\tok{e}\sep\tok{s}\sep\tok{o}\to\tok{p}\sep\tok{e}\sep\tok{so}$. «sopa» $=\tok{s}\sep\tok{o}\sep\tok{p}\sep\tok{a}\to\tok{so}\sep\tok{p}\sep\tok{a}\to\tok{so}\sep\tok{pa}$. «oso» $=\tok{o}\sep\tok{s}\sep\tok{o}\to\tok{o}\sep\tok{so}$.
\end{enumerate}
\resultado \textbf{fusiones $\tok{s}+\tok{o}$ (12), $\tok{p}+\tok{i}$ (6), $\tok{pi}+\tok{so}$ (6) y $\tok{p}+\tok{a}$ (6); símbolos $53\to41\to35\to29\to23$; «peso» $\to\tok{p}\sep{}\tok{e}\sep{}\tok{so}$, «sopa» $\to\tok{so}\sep{}\tok{pa}$, «oso» $\to\tok{o}\sep{}\tok{so}$.}
\comprobacion recuento del corpus final, \tok{pa}\sep\tok{so}$\times4$, \tok{pa}\sep\tok{so}\sep\tok{s}$\times2$, \tok{piso}$\times3$, \tok{piso}\sep\tok{s}$\times3$: $8+6+3+6=23$. Con la regla «el primer par que aparece», el paso~2 habría elegido $(\tok{p},\tok{a})$ y el tokenizador sería otro.
\end{solucion}

\begin{solucion}{eje:2-compresion}
\datos secuencia de longitud $L$; par $(a,b)$ con $a\neq b$ y $c$ apariciones en las posiciones $i_1<\dots<i_c$ (cada $i_r$ es la posición de la $a$).
\planteamiento probar que $i_{r+1}\ge i_r+2$, de modo que las sustituciones no comparten símbolos, y contar.
\pasos
\begin{enumerate}[label=\arabic*.]
  \item Si $i_{r+1}=i_r+1$, el símbolo de la posición $i_r+1$ sería $b$ (segundo del par $r$) y $a$ (primero del par $r+1$), luego $a=b$: contradicción. Así $i_{r+1}\ge i_r+2$.
  \item Los intervalos $\{i_r,i_r+1\}$ son, por tanto, disjuntos, y la fusión de izquierda a derecha sustituye los $c$.
  \item Cada sustitución cambia 2 símbolos por 1 y resta 1; en total, $L'=L-c$.
  \item Contraejemplo con $a=b$: «\tok{aaa}» tiene $c=2$ pares $(\tok{a},\tok{a})$ que se solapan; la fusión da \tok{aa}\sep\tok{a}, de longitud $2=3-1\neq 3-2$.
  \item «\tok{aaaa}» tiene 3 pares; la fusión da \tok{aa}\sep\tok{aa}, de longitud $2=4-2$. «\tok{aaaaa}» tiene 4 pares; la fusión da \tok{aa}\sep\tok{aa}\sep\tok{a}, de longitud $3=5-2$.
\end{enumerate}
\resultado \textbf{$L'=L-c$ si $a\neq b$}; con $a=b$ solo cuentan las apariciones no solapadas, $\lfloor r/2\rfloor$ en cada racha de $r$ símbolos iguales: longitudes 2 y 3 para «\tok{aaaa}» y «\tok{aaaaa}».
\comprobacion en el \cref{tab:bpe-casa}, $30-8=22$, $22-5=17$, $17-3=14$ y $14-2=12$ coinciden con el recuento directo.
\end{solucion}

\begin{solucion}{eje:2-medio}
\datos tokenizador BPE de bytes con el token \hx{C3}.
\planteamiento usar el \cref{tab:utf8} para (a), el \cref{teo:bpe-correcto} y el \cref{teo:prefijo} para (b).
\pasos
\begin{enumerate}[label=\arabic*.]
  \item (a) $\hx{C3}=\bits{11000011}$ anuncia un carácter de 2 bytes; aislado, falta el byte de continuación y la cadena no es UTF-8 válida. La decodificación estricta da error y la tolerante devuelve el carácter de sustitución \hx{U+FFFD}.
  \item (b) Sea $b=\Enc(x)$ la cadena UTF-8 del texto $x$ y sea $(t_1,\dots,t_m)$ su tokenización. Por el \cref{teo:bpe-correcto} con $\V_0=\mathcal{B}$, $t_1\Vert\dots\Vert t_m=b$.
  \item Por el \cref{teo:prefijo}, $\Enc$ es inyectiva y $b$ es una cadena válida, luego $\Enc^{-1}(t_1\Vert\dots\Vert t_m)=\Enc^{-1}(b)=x$.
  \item (c) Si la «ñ» $=\hx{C3 B1}$ llega en dos tokens y cada uno se decodifica al recibirlo, se muestran dos caracteres \hx{U+FFFD} en lugar de «ñ».
  \item Remedio: acumular en un búfer los bytes de un carácter incompleto (se sabe que lo está por su byte inicial, paso~1 de la demostración del \cref{teo:prefijo}) y mostrarlo cuando llegue su último byte.
\end{enumerate}
\resultado \textbf{(a) error o \hx{U+FFFD}; (b) la concatenación de los tokens es la cadena UTF-8 original, que se decodifica sin pérdida; (c) el usuario ve dos \hx{U+FFFD}; se corrige con un búfer de bytes.}
\comprobacion en Python, decodificar por separado los dos bytes de \texttt{"ñ".encode()} con \texttt{errors="replace"} da dos \hx{U+FFFD}, mientras que \texttt{"ñ".encode().decode()} devuelve «ñ».
\end{solucion}

% ---------------------------------------------------------------------
\section{Tema 3}

\begin{solucion}{eje:3-palabras}
\datos un LLM multilingüe con vocabulario de palabras completas.
\planteamiento comparar lo que exige ese vocabulario con lo que necesita el modelo: un $K$ manejable y capacidad para representar cualquier texto.
\pasos
\begin{enumerate}[label=\arabic*.]
  \item Tamaño: cada forma flexionada («canto», «cantas», «cantaré»\dots) es un token, y con decenas de idiomas $K$ supera con creces el medio millón. Por la \cref{prop:coste-vocab}, las matrices de entrada y salida crecen en proporción.
  \item Cobertura: una palabra nueva, una errata o un nombre propio raro no estarían en $\V$; habría que sustituirlos por un token «desconocido» y se perdería su información.
  \item Además, «caballería» y «caballerías» tendrían filas independientes en $E$, y las formas raras recibirían muy pocos ejemplos de entrenamiento.
\end{enumerate}
\resultado \textbf{vocabulario gigantesco (matrices de embeddings y de salida enormes) y palabras fuera del vocabulario sin representación}; además, las formas emparentadas no comparten parámetros.
\comprobacion las subpalabras resuelven los tres problemas: $K$ entre $3{,}2\cdot 10^{4}$ y $2{,}6\cdot 10^{5}$ y toda palabra se descompone en piezas conocidas (\cref{cor:no-unk}).
\end{solucion}

\begin{solucion}{eje:3-letras}
\datos «ferrocarril» tiene 11 letras, 4 de ellas «r».
\planteamiento el modelo recibe tokens, no letras; hay que ver qué información falta en la entrada.
\pasos
\begin{enumerate}[label=\arabic*.]
  \item El tokenizador divide la palabra en pocas piezas, por ejemplo \tok{ferro}\,\sep{}\,\tok{carril}, y el modelo recibe dos enteros.
  \item Cada entero se convierte en una fila de $E$, que no contiene una lista explícita de letras.
  \item Para contar, el modelo debe saber que \tok{ferro} contiene 2 «r» y \tok{carril} otras 2, y sumarlas. Esa ortografía solo se aprende de forma indirecta, de textos que deletrean palabras.
  \item Remedio: escribir la palabra separada letra a letra; entonces cada letra es, casi siempre, un token.
\end{enumerate}
\resultado \textbf{falla porque no ve letras sino tokens; contar exige conocer la ortografía interna de cada token, que no está en la entrada.}
\comprobacion con la palabra deletreada, el número de tokens que son «r» coincide con la respuesta, $4$.
\end{solucion}

\begin{solucion}{eje:3-vf}
\planteamiento contrastar cada afirmación con la \cref{def:cod-bpe}, la \cref{obs:congelado} y la \cref{prop:coste-vocab}.
\pasos
\begin{enumerate}[label=(\alph*)]
  \item \textbf{Falsa.} BPE aplica las fusiones en un orden fijo y es determinista; el modelo unigram es el que admite varias segmentaciones y permite muestrearlas.
  \item \textbf{Falsa.} El tokenizador se entrena antes, con un algoritmo de conteo, y se congela.
  \item \textbf{Verdadera.} La matriz tiene $Kd$ parámetros; con $d$ fijo, $(2K)d=2(Kd)$.
\end{enumerate}
\resultado \textbf{(a) falsa; (b) falsa; (c) verdadera.}
\comprobacion en (c), con $d=4096$: $32\,000\cdot 4096=131\,072\,000$ y $64\,000\cdot 4096=262\,144\,000$, exactamente el doble.
\end{solucion}

\begin{solucion}{eje:3-perdida}
\datos vocabulario del \cref{ej:unigram}; $s^{*}=(\tok{casa},\tok{s})$.
\planteamiento $-\ln P(s)=-\sum_i\ln p(t_i)$.
\pasos
\begin{enumerate}[label=\arabic*.]
  \item $P(s^{*})=0{,}2\cdot 0{,}1=0{,}02$ y $-\ln 0{,}02=\ln 100-\ln 2=4{,}6052-0{,}6931=3{,}912$ nats.
  \item $P(\tok{cas},\tok{as})=0{,}1\cdot 0{,}15=0{,}015$ y $-\ln 0{,}015=\ln 100-\ln 1{,}5=4{,}6052-0{,}4055=4{,}200$ nats.
  \item Diferencia: $\ln(0{,}02/0{,}015)=\ln(4/3)=0{,}2877$ nats.
\end{enumerate}
\resultado \textbf{$-\ln P(s^{*})=3{,}912$ nats; con $(\tok{cas},\tok{as})$ se pierden $\ln(4/3)\approx 0{,}2877$ nats.}
\comprobacion la diferencia es positiva porque $s^{*}$ es la segmentación de máxima probabilidad. Restar valores ya redondeados, $4{,}2-3{,}9=0{,}3$, daría una cifra demasiado gruesa.
\end{solucion}

\begin{solucion}{eje:3-param}
\datos $K=32\,000$, $d=4096$, matrices separadas, 2 bytes por parámetro.
\planteamiento cada matriz tiene $Kd$ parámetros (\cref{prop:coste-vocab}).
\pasos
\begin{enumerate}[label=\arabic*.]
  \item Una matriz: $32\,000\cdot 4096=131\,072\,000$.
  \item Dos matrices: $2\cdot 131\,072\,000=262\,144\,000$.
  \item Memoria: $262\,144\,000\cdot 2=524\,288\,000$ bytes $=0{,}5243$~GB ($10^{9}$ bytes), o $0{,}4883$~GiB ($2^{30}$ bytes).
\end{enumerate}
\resultado \textbf{$262\,144\,000$ parámetros, que ocupan unos $0{,}52$~GB en \texttt{bf16}.}
\comprobacion coincide con la fila de Llama~2 7B del \cref{tab:vocab-modelos}: $131$ millones por matriz.
\end{solucion}

\begin{solucion}{eje:3-numeros}
\datos agrupación de los dígitos en bloques de hasta 3, de izquierda a derecha.
\planteamiento aplicar la regla a cada número y observar a qué posición decimal corresponde cada dígito dentro de su token.
\pasos
\begin{enumerate}[label=\arabic*.]
  \item (a) $1605\to\tok{160}\,\sep{}\,\tok{5}$; $2026\to\tok{202}\,\sep{}\,\tok{6}$; $123456\to\tok{123}\,\sep{}\,\tok{456}$; $1234567\to\tok{123}\,\sep{}\,\tok{456}\,\sep{}\,\tok{7}$.
  \item (b) $999\to\tok{999}$, un token; $999+1=1000\to\tok{100}\,\sep{}\,\tok{0}$, dos tokens que no tienen ningún bloque en común con el anterior.
  \item (c) Con la agrupación de izquierda a derecha, el último token de $1605$ son las unidades, pero el último de $123456$ son centenas, decenas y unidades: la posición decimal de un dígito depende de la longitud del número, y el acarreo de una suma puede cambiar la segmentación entera.
  \item Una agrupación de derecha a izquierda, como la de los separadores de millares, alinea los bloques con las potencias de $10^{3}$: $1605\to\tok{1}\,\sep{}\,\tok{605}$, $1234567\to\tok{1}\,\sep{}\,\tok{234}\,\sep{}\,\tok{567}$.
\end{enumerate}
\resultado \textbf{(a) \tok{160}\sep{}\tok{5}, \tok{202}\sep{}\tok{6}, \tok{123}\sep{}\tok{456}, \tok{123}\sep{}\tok{456}\sep{}\tok{7}; (b) 1 token frente a 2; (c) la agrupación de izquierda a derecha desalinea las posiciones decimales; la de derecha a izquierda las respeta.}
\comprobacion con la agrupación de derecha a izquierda, $999\to\tok{999}$ y $1000\to\tok{1}\,\sep{}\,\tok{000}$: el bloque de unidades, decenas y centenas ocupa siempre el último token.
\end{solucion}

\begin{solucion}{eje:3-viterbi}
\datos cadena de $n$ caracteres; vocabulario con todos los caracteres y tokens de longitud a lo sumo $\ell$.
\planteamiento (a) biyección entre segmentaciones y conjuntos de cortes; (b) asociar a cada segmentación la composición de $n$ dada por las longitudes de sus tokens.
\pasos
\begin{enumerate}[label=\arabic*.]
  \item (a) Una segmentación en subcadenas no vacías queda determinada por el conjunto de los huecos entre caracteres consecutivos en los que se corta. Hay $n-1$ huecos y cada uno se corta o no, de forma independiente: $2^{n-1}$ segmentaciones.
  \item (b) A cada segmentación con tokens de $\V$ se le asocia la secuencia de longitudes $(\ell_1,\dots,\ell_m)$, con $\ell_i\in\{1,\dots,\ell\}$ y $\sum\ell_i=n$. La aplicación es inyectiva, porque las longitudes determinan los puntos de corte; luego el número de segmentaciones es a lo sumo el número $c_{\ell}(n)$ de composiciones de $n$ con partes en $\{1,\dots,\ell\}$.
  \item $c_{\ell}$ satisface $c_{\ell}(0)=1$ y $c_{\ell}(n)=\sum_{j=1}^{\min(\ell,n)}c_{\ell}(n-j)$, porque la última parte vale $j$. Con $\ell=2$ es la sucesión de Fibonacci, $c_2(n)=F_{n+1}$, que crece exponencialmente; por ejemplo, $c_2(20)=F_{21}=10\,946$.
  \item Si $\V$ contiene todas las subcadenas de longitud a lo sumo $\ell$, la cota se alcanza, y enumerar las segmentaciones cuesta un tiempo exponencial en $n$. El \cref{teo:viterbi} encuentra la mejor con $O(n\ell)$ productos.
\end{enumerate}
\resultado \textbf{(a) $2^{n-1}$; (b) a lo sumo $c_{\ell}(n)$, que crece exponencialmente, frente al coste lineal $O(n\ell)$ de Viterbi.}
\comprobacion con $n=5$ y $\ell=5$, $c_5(5)=2^{4}=16$ composiciones; con el vocabulario del \cref{ej:unigram}, solo 7 de las 16 segmentaciones de «casas» usan tokens de $\V$ (\cref{tab:unigram}).
\end{solucion}

% ---------------------------------------------------------------------
\section{Tema 4}

\begin{solucion}{eje:4-onehot}
\datos $E$ de tamaño $3\times 2$; $e_1=(0,1,0)^{\top}$.
\planteamiento calcular el producto componente a componente y contar las multiplicaciones de un producto $e_i^{\top}E$ general.
\pasos
\begin{enumerate}[label=\arabic*.]
  \item Primera componente: $0\cdot 1+1\cdot 2+0\cdot 0=2$. Segunda: $0\cdot 0+1\cdot(-1)+0\cdot 3=-1$.
  \item Luego $e_1^{\top}E=(2,-1)=E_{1,:}$, la fila de índice~1.
  \item El producto general requiere $K$ multiplicaciones por cada una de las $d$ componentes: $Kd=50\,257\cdot 768=38\,597\,376$.
\end{enumerate}
\resultado \textbf{$e_1^{\top}E=(2,-1)$; el producto completo costaría $38\,597\,376$ multiplicaciones}, frente a ninguna con la indexación.
\comprobacion la cifra coincide con el número de parámetros de la matriz de GPT-2 pequeño (\cref{tab:vocab-modelos}), como debe ser: hay una multiplicación por entrada de $E$.
\end{solucion}

\begin{solucion}{eje:4-pares}
\datos «la luna brilla», $n=3$; ventanas $m=1$ y $m=2$.
\planteamiento emparejar cada palabra con las que están a distancia a lo sumo $m$.
\pasos
\begin{enumerate}[label=\arabic*.]
  \item $m=1$: $(\text{la},\text{luna})$; $(\text{luna},\text{la})$, $(\text{luna},\text{brilla})$; $(\text{brilla},\text{luna})$. Total: $1+2+1=4$.
  \item $m=2$: se añaden $(\text{la},\text{brilla})$ y $(\text{brilla},\text{la})$. Total: $6$.
  \item \cref{prop:pares}: $1\cdot(6-1-1)=4$ y $2\cdot(6-2-1)=6$; se cumple $n\ge m+1$ en ambos casos.
\end{enumerate}
\resultado \textbf{4 pares con ventana~1 y 6 con ventana~2.}
\comprobacion con $m=2$ cada palabra se empareja con las otras dos: $3\cdot 2=6$ pares ordenados.
\end{solucion}

\begin{solucion}{eje:4-coseno}
\datos cereza $=(40,3)$, fresa $=(20,4)$, digital $=(2,30)$.
\planteamiento $\cos(a,b)=a\cdot b/(\norm{a}\norm{b})$.
\pasos
\begin{enumerate}[label=\arabic*.]
  \item Normas: $\norm{\text{cereza}}=\sqrt{1609}=40{,}11$; $\norm{\text{fresa}}=\sqrt{416}=20{,}40$; $\norm{\text{digital}}=\sqrt{904}=30{,}07$.
  \item $\text{cereza}\cdot\text{fresa}=800+12=812$ y $\sqrt{1609\cdot 416}=\sqrt{669\,344}=818{,}1$; coseno $812/818{,}1=0{,}9925$.
  \item $\text{cereza}\cdot\text{digital}=80+90=170$ y $\sqrt{1609\cdot 904}=\sqrt{1\,454\,536}=1206$; coseno $170/1206=0{,}1410$.
  \item (b) Para $\lambda>0$: $(\lambda a)\cdot b=\lambda(a\cdot b)$ y $\norm{\lambda a}=\abs{\lambda}\norm{a}=\lambda\norm{a}$; luego $\cos(\lambda a,b)=\lambda(a\cdot b)/(\lambda\norm{a}\norm{b})=\cos(a,b)$.
  \item La propiedad es deseable porque la longitud de un vector de recuentos crece con la frecuencia de la palabra, y la similitud no debe depender de cuán frecuente sea cada palabra.
\end{enumerate}
\resultado \textbf{$\cos(\text{cereza},\text{fresa})\approx 0{,}9925$ y $\cos(\text{cereza},\text{digital})\approx 0{,}1410$: cereza se parece mucho más a fresa.}
\comprobacion $3\cdot\text{cereza}=(120,9)$ da el mismo coseno con fresa: $(2400+36)/\sqrt{14\,481\cdot 416}=2436/2454=0{,}9925$.
\end{solucion}

\begin{solucion}{eje:4-pmi}
\datos $N(\text{gato},\text{come})=N(\text{gato},\text{maúlla})=N(\text{perro},\text{come})=N(\text{perro},\text{ladra})=4$; el resto, 0; $k=2$.
\planteamiento calcular las probabilidades conjuntas y marginales y aplicar la \cref{def:pmi} y el \cref{teo:levy}.
\pasos
\begin{enumerate}[label=\arabic*.]
  \item Total: $\abs{\mathcal{D}}=16$. Marginales: $P(\text{gato})=P(\text{perro})=8/16=1/2$; $P(\text{come})=8/16=1/2$; $P(\text{maúlla})=P(\text{ladra})=4/16=1/4$.
  \item $\PMI(\text{gato},\text{come})=\ln\dfrac{4/16}{(1/2)(1/2)}=\ln 1=0$.
  \item $\PMI(\text{gato},\text{maúlla})=\ln\dfrac{4/16}{(1/2)(1/4)}=\ln 2=0{,}6931$.
  \item $\PMI(\text{gato},\text{ladra})=\ln 0=-\infty$, luego $\PPMI=0$. Por simetría, «perro» da $0$ con «come», $\ln 2$ con «ladra» y $-\infty$ con «maúlla».
  \item (b) Para los pares observados, $u_c\cdot v_w=\PMI-\ln 2$: $0-0{,}6931=-0{,}6931$ con «come» y $0{,}6931-0{,}6931=0$ con «maúlla» o «ladra».
\end{enumerate}
\resultado \textbf{$\PMI=0$ con «come», $\ln 2\approx 0{,}6931$ con el verbo propio de cada animal y $-\infty$ ($\PPMI=0$) con el ajeno; SGNS con $k=2$ aprendería productos $-0{,}6931$ y $0$.}
\comprobacion «come» acompaña a ambos animales con la misma frecuencia, de modo que no informa: su PMI es nula, mientras que «maúlla» duplica la frecuencia esperada con «gato» ($e^{\ln 2}=2$).
\end{solucion}

\begin{solucion}{eje:4-softmax}
\datos $J=-\ln\hat y_o$ con $\hat y=\softmax(U^{\top}v_c)$; $u_{\text{come}}=(0,1)$, $u_{\text{avión}}=(-1,0)$, $v_c=(1,0)$, $o=\text{come}$.
\planteamiento partir de la \cref{prop:grad-vc} y derivar $J=-u_o^{\top}v_c+\ln\sum_w\exp(u_w^{\top}v_c)$ respecto a cada $u_w$.
\pasos
\begin{enumerate}[label=\arabic*.]
  \item (a) $\sum_w\hat y_w u_w=U\hat y$ y $u_o=Uy$, luego $\partial J/\partial v_c=U(\hat y-y)$. Se anula si y solo si $u_o=\sum_w\hat y_w u_w$, es decir, si el vector de contexto real coincide con la media de los vectores de contexto ponderada por la predicción.
  \item (b) El término $-u_o^{\top}v_c$ solo depende de $u_o$ y da $-v_c$; el segundo da $\hat y_w v_c$. Luego $\partial J/\partial u_o=(\hat y_o-1)v_c$ y $\partial J/\partial u_w=\hat y_w v_c$ si $w\neq o$; en una sola fórmula, $(\hat y_w-y_w)v_c$.
  \item (c) Puntuaciones: $u_{\text{come}}\cdot v_c=0$ y $u_{\text{avión}}\cdot v_c=-1$. Softmax: $\hat y=(e^{0},e^{-1})/(1+e^{-1})=(0{,}7311;\ 0{,}2689)$.
  \item $\partial J/\partial v_c=0{,}7311\,(0,1)+0{,}2689\,(-1,0)-(0,1)=(-0{,}2689;\ -0{,}2689)$.
  \item $\partial J/\partial u_{\text{come}}=(0{,}7311-1)(1,0)=(-0{,}2689;\ 0)$ y $\partial J/\partial u_{\text{avión}}=0{,}2689\,(1,0)=(0{,}2689;\ 0)$.
\end{enumerate}
\resultado \textbf{$\partial J/\partial v_c=U(\hat y-y)=(-0{,}2689;\,-0{,}2689)$; $\partial J/\partial u_w=(\hat y_w-y_w)v_c$, que vale $(-0{,}2689;\,0)$ para «come» y $(0{,}2689;\,0)$ para «avión».}
\comprobacion la pérdida inicial es $-\ln 0{,}7311=0{,}3133$. Un paso con $\eta=0{,}1$ sobre $v_c$ da $v_c'=(1{,}0269;\ 0{,}0269)$, puntuaciones $0{,}0269$ y $-1{,}0269$, y pérdida $0{,}2991<0{,}3133$.
\end{solucion}

\begin{solucion}{eje:4-sgns}
\datos $L=-\ln\sigma(c_{+}\cdot w)-\ln\sigma(-c_{-}\cdot w)$; $w=(1,0)$, $c_{+}=(0,1)$, $c_{-}=(-1,0)$; $\eta=0{,}1$.
\planteamiento regla de la cadena con $\sigma'(x)=\sigma(x)(1-\sigma(x))$, de donde $\frac{d}{dx}[-\ln\sigma(x)]=\sigma(x)-1$, y $1-\sigma(-x)=\sigma(x)$.
\pasos
\begin{enumerate}[label=\arabic*.]
  \item Término positivo, con $s_{+}=c_{+}\cdot w$: su derivada respecto a $s_{+}$ es $\sigma(s_{+})-1$, y $\partial s_{+}/\partial c_{+}=w$. Luego $\partial L/\partial c_{+}=(\sigma(c_{+}\cdot w)-1)\,w$.
  \item Término negativo, con $s_{-}=-c_{-}\cdot w$: su derivada es $\sigma(s_{-})-1=-\sigma(c_{-}\cdot w)$, y $\partial s_{-}/\partial c_{-}=-w$. Luego $\partial L/\partial c_{-}=\sigma(c_{-}\cdot w)\,w$.
  \item Respecto a $w$ intervienen ambos términos: $\partial L/\partial w=(\sigma(c_{+}\cdot w)-1)\,c_{+}+\sigma(c_{-}\cdot w)\,c_{-}$.
  \item (b) $c_{+}\cdot w=0$ y $c_{-}\cdot w=-1$; $\sigma(0)=0{,}5$ y $\sigma(-1)=1/(1+e)=0{,}2689$.
  \item $\partial L/\partial c_{+}=-0{,}5\,(1,0)=(-0{,}5;\ 0)$; $\partial L/\partial c_{-}=0{,}2689\,(1,0)=(0{,}2689;\ 0)$; $\partial L/\partial w=-0{,}5\,(0,1)+0{,}2689\,(-1,0)=(-0{,}2689;\ -0{,}5)$.
  \item (c) $w'=w-\eta\,\partial L/\partial w=(1+0{,}02689;\ 0+0{,}05)=(1{,}027;\ 0{,}05)$.
  \item Nueva pérdida: $c_{+}\cdot w'=0{,}05$ da $-\ln\sigma(0{,}05)=0{,}6685$; $c_{-}\cdot w'=-1{,}027$ da $-\ln\sigma(1{,}027)=0{,}3061$. $L'=0{,}9746$.
\end{enumerate}
\resultado \textbf{$\partial L/\partial c_{+}=(-0{,}5;\,0)$, $\partial L/\partial c_{-}=(0{,}2689;\,0)$, $\partial L/\partial w=(-0{,}2689;\,-0{,}5)$; tras un paso, $L$ baja de $1{,}006$ a $0{,}9746$.}
\comprobacion el paso acerca $w$ a $c_{+}$ (su segunda componente sube de 0 a $0{,}05$) y lo aleja de $c_{-}$, que apunta en el sentido negativo del primer eje. Un gradiente numérico por diferencias centradas da los mismos vectores.
\end{solucion}

% ---------------------------------------------------------------------
\section{Tema 5}

\begin{solucion}{eje:5-dispersion}
\datos $N=440$ palabras, $K=243$ tipos, $416$ bigramas distintos, $398$ únicos.
\planteamiento aplicar la \cref{prop:dispersion} con $n=2$ y recordar el EMV del módulo~1.
\pasos
\begin{enumerate}[label=\arabic*.]
  \item (a) Cota: $N-n+1=439$ bigramas distintos como máximo, y se observan $416\le 439$. Fracción máxima: $439/243^{2}=439/59\,049=0{,}743\,\%$; la observada es $416/59\,049=0{,}704\,\%$.
  \item (b) $398/416=0{,}9567$: el $95{,}67\,\%$ de los bigramas observados aparece una sola vez.
  \item (c) El EMV asigna probabilidad~0 a todo bigrama no observado. Si el texto de evaluación contiene uno, su probabilidad es~0, la pérdida es infinita y la perplejidad también (módulo~1).
\end{enumerate}
\resultado \textbf{(a) $416\le 439$ y $0{,}704\,\%\le 0{,}743\,\%$; (b) $95{,}67\,\%$; (c) probabilidad nula y perplejidad infinita.}
\comprobacion casi todos los recuentos positivos valen~1, de modo que la estimación de las pocas casillas no vacías tampoco es fiable: la dispersión no se arregla solo con suavizado.
\end{solucion}

\begin{solucion}{eje:5-clip}
\datos $g=(3,-4,12)$, $\tau=5$.
\planteamiento calcular $\norm{g}$ y, como supera $\tau$, reescalar.
\pasos
\begin{enumerate}[label=\arabic*.]
  \item $\norm{g}=\sqrt{9+16+144}=\sqrt{169}=13\ge 5$.
  \item $\clip_5(g)=\frac{5}{13}(3,-4,12)=(15/13,\,-20/13,\,60/13)=(1{,}154;\ -1{,}538;\ 4{,}615)$.
  \item Norma: $\frac{5}{13}\cdot 13=5$.
  \item El vector es $\lambda g$ con $\lambda=5/13>0$: misma dirección y sentido.
\end{enumerate}
\resultado \textbf{$\clip_5(g)=(15/13,\,-20/13,\,60/13)\approx(1{,}154;\,-1{,}538;\,4{,}615)$, de norma~5.}
\comprobacion $(15^{2}+20^{2}+60^{2})/13^{2}=(225+400+3600)/169=4225/169=25$, y $\sqrt{25}=5$.
\end{solucion}

\begin{solucion}{eje:5-olvido}
\datos $f$ constante; 100 pasos; RNN escalar con $w=0{,}9$.
\planteamiento por la \cref{prop:lstm-grad}, el factor del camino directo es $f^{100}$.
\pasos
\begin{enumerate}[label=\arabic*.]
  \item (a) $0{,}98^{100}=e^{100\ln 0{,}98}=e^{-2{,}020}=0{,}1326$. En la RNN, $0{,}9^{100}=2{,}656\cdot 10^{-5}$ (\cref{tab:wT}).
  \item Cociente: $0{,}1326/2{,}656\cdot 10^{-5}=4993\approx 5{,}0\cdot 10^{3}$.
  \item (b) $f^{100}\ge 0{,}5\iff 100\ln f\ge\ln 0{,}5=-0{,}6931\iff f\ge e^{-0{,}006931}=0{,}993092$.
  \item Para dar un valor seguro con 4 cifras decimales se redondea hacia arriba: $f\ge 0{,}9931$.
\end{enumerate}
\resultado \textbf{(a) $0{,}1326$, unas $5\cdot 10^{3}$ veces más que en la RNN; (b) $f\ge 0{,}5^{1/100}\approx 0{,}99309$.}
\comprobacion $0{,}9931^{100}=0{,}5004\ge 0{,}5$, mientras que $0{,}9930^{100}=0{,}4954<0{,}5$.
\end{solucion}

\begin{solucion}{eje:5-param}
\datos $h=d=256$; la GRU tiene tres transformaciones como las de la RNN.
\planteamiento \cref{prop:param-rnn} y \cref{cor:param-lstm}.
\pasos
\begin{enumerate}[label=\arabic*.]
  \item RNN: $h(h+d)+h=256\cdot 512+256=131\,072+256=131\,328$.
  \item LSTM: $4\cdot 131\,328=525\,312$.
  \item GRU: $3\cdot 131\,328=393\,984$.
  \item Cociente LSTM/RNN: $525\,312/131\,328=4$.
\end{enumerate}
\resultado \textbf{RNN $131\,328$, GRU $393\,984$ y LSTM $525\,312$ parámetros; un paso de LSTM cuesta unas 4 veces más que uno de RNN.}
\comprobacion los tres modelos siguen siendo secuenciales: el factor 4 aumenta el coste de cada paso, pero no cambia el número de pasos, $T$.
\end{solucion}

\begin{solucion}{eje:5-wmin}
\datos $h_t=w\,h_{t-1}$, $0<w<1$; factor a 100 pasos $w^{100}\ge 0{,}1$.
\planteamiento tomar logaritmos; $\ln$ es creciente y conserva la desigualdad.
\pasos
\begin{enumerate}[label=\arabic*.]
  \item $100\ln w\ge\ln 0{,}1=-2{,}3026$.
  \item $\ln w\ge -0{,}023026$, luego $w\ge e^{-0{,}023026}=0{,}977237$.
  \item Con 4 cifras decimales, redondear hacia abajo no basta: $0{,}9772^{100}=0{,}0996<0{,}1$, mientras que $0{,}9773^{100}=0{,}1006$. El valor seguro es $0{,}9773$.
  \item Interpretación: con $w<0{,}9772$ llega menos del $10\,\%$ de la señal a 100 pasos, y con $w>1$ el factor explota ($1{,}1^{100}\approx 13\,781$). La franja útil, entre $0{,}977$ y~$1$, es estrechísima.
\end{enumerate}
\resultado \textbf{$w\ge 0{,}1^{1/100}\approx 0{,}97724$} (con 4 cifras decimales, redondeando hacia arriba, $0{,}9773$).
\comprobacion $0{,}97724^{100}=e^{100\ln 0{,}97724}=e^{-2{,}3023}=0{,}1000$. Las LSTM y las conexiones residuales fijan ese factor cerca de~1 por diseño (\cref{prop:lstm-grad}).
\end{solucion}

\begin{solucion}{eje:5-autovalores}
\datos $M$ simétrica con $Mv_i=\lambda_i v_i$, $\{v_i\}$ ortonormal, $\abs{\lambda_1}>\abs{\lambda_2}\ge\dots$; $x=\sum_i a_i v_i$ con $a_1\neq 0$.
\planteamiento (a) inducción sobre $k$; (b) descomponer $M^{k}x$ en la base de vectores propios; (c) aplicar el resultado al jacobiano $\partial h_T/\partial h_0=M^{T}$.
\pasos
\begin{enumerate}[label=\arabic*.]
  \item (a) $M^{1}v_i=\lambda_i v_i$. Si $M^{k-1}v_i=\lambda_i^{k-1}v_i$, entonces $M^{k}v_i=M(\lambda_i^{k-1}v_i)=\lambda_i^{k-1}\lambda_i v_i=\lambda_i^{k}v_i$. Como los $v_i$ forman una base, estos son todos los valores propios de $M^{k}$.
  \item (b) Por linealidad, $M^{k}x=a_1\lambda_1^{k}v_1+r_k$ con $r_k=\sum_{i\ge 2}a_i\lambda_i^{k}v_i$, que es ortogonal a $v_1$.
  \item La tangente del ángulo $\theta_k$ entre $M^{k}x$ y $v_1$ es $\norm{r_k}/\abs{a_1\lambda_1^{k}}$. Por la ortonormalidad, $\norm{r_k}^{2}=\sum_{i\ge 2}a_i^{2}\lambda_i^{2k}\le\abs{\lambda_2}^{2k}\sum_{i\ge 2}a_i^{2}=\abs{\lambda_2}^{2k}\norm{x-a_1v_1}^{2}$.
  \item Luego $\tan\theta_k\le(\abs{\lambda_2}/\abs{\lambda_1})^{k}\,\norm{x-a_1v_1}/\abs{a_1}$, que tiende a~0. Para un $x$ aleatorio con distribución continua, $a_1\neq 0$ con probabilidad~1.
  \item (c) En la RNN lineal $h_t=Mh_{t-1}$, $\partial h_T/\partial h_0=M^{T}$, y el gradiente que llega a $h_0$ es $M^{T}g$ con $g=\partial L/\partial h_T$. Por (b), se alinea con $v_1$ y su norma se comporta como $\abs{a_1}\abs{\lambda_1}^{T}$: se desvanece si $\abs{\lambda_1}<1$ y explota si $\abs{\lambda_1}>1$.
\end{enumerate}
\resultado \textbf{(a) $M^{k}v_i=\lambda_i^{k}v_i$; (b) $\tan\theta_k\le(\abs{\lambda_2}/\abs{\lambda_1})^{k}\norm{x-a_1v_1}/\abs{a_1}\to 0$; (c) el gradiente de una RNN lineal decae o crece como $\abs{\lambda_1}^{T}$.}
\comprobacion para una matriz simétrica la norma espectral es $\abs{\lambda_1}$, de modo que (c) es coherente con la cota $\norm{W}^{T}$ del \cref{teo:desvanece}. Con $M=\diag(0{,}9;\,0{,}5)$ y $x=(1,1)$, $M^{10}x=(0{,}3487;\ 0{,}000977)$ y $\tan\theta_{10}=0{,}0028=(0{,}5/0{,}9)^{10}$.
\end{solucion}
